
This project keeps a correction register. Earlier versions of this work
carried it in the main text; a referee for this journal observed that
scientific transparency should be visible in the design and the supplement
rather than repeatedly asserted, and that is right, so it is here.

\texttt{REFEREE-LOG.md} in the archive carries every objection this project
has received, each with its disposition and with the corrections it produced
named. What follows is the register grouped by \emph{class}, because the
classes are more useful than the incidents: they say what kind of mistake this
apparatus is prone to. Identifiers run in the order the corrections were made,
across all twenty rounds; three of the low numbers never entered the register
and the classes below are therefore not a contiguous range, which is a fact
about the numbering and not a missing entry.

\paragraph{Class A --- a control that could not fail}
Four corrections (C1, C2, C8, C16) are nulls or searches drawn at the wrong
level, swept over the wrong range, or given fresh noise inside a comparison
that assumed determinism. A null that cannot fail is worse than no null,
because it is reported as evidence. The most recent, C16, is in
\ref{app:props}: the axis non-redundancy search drew fresh noise
inside each cell evaluation, so two evaluations of the same cell disagreed and
the search found almost nothing.

\paragraph{Class B --- a quantity whose sign was set by a tie-break}
Two corrections (C4, C7) concern an operational factor computed at a fixed
review capacity, where $93.1\%$ of what the naive baseline nominated came from
a single tied block and reordering rows inside that block moved the arm from
$-26$ to $+608$. A quantity whose sign is set by a tie-break is not a
measurement. Section~\ref{sec:dca} is built on net benefit, which admits or
excludes a whole tied block at a threshold and never breaks one.

\paragraph{Class C --- every literal correct, the relation between them
false}
Five corrections (C9, C10, C12, C13, C14) are sentences in which each printed
number was right and the claim asserted about them was not: an extremum that
was the worst of the points the table happened to name; a count reported as a
contiguous run; a range ``across the operating range'' that was a range across
five named points; a universal ``if and only if'' whose necessity direction is
only existential (Section~\ref{sec:props}); and --- the one that mattered most
--- a headline spread that was inflated by an intercept-only baseline no
analyst would build, while a later section of the same manuscript said so and
the abstract did not. This class is the reason
Section~\ref{sec:regionsdef} excludes the intercept-only rung by definition and
the reason every number in this document is a generated macro rather than a
typed literal.

\paragraph{Class D --- evidence one download away}
One correction (C11): an earlier version reported that it could not establish
whether the knowledge-article reference was available at incident creation,
named the file that would settle where the field comes from, and did not
obtain it. Obtaining it changed the headline. Section~\ref{sec:cmdb} now
reports two decision times rather than choosing one, because obtaining the
file settled where the field is written and not when.

\paragraph{Class E --- a claim about the world that our own coder could not
support}
One correction (C15), and it is the largest. An earlier version reported that
\emph{not one} of twenty audited papers reports the increment across a range
of operating points. The rebuilt pilot of Section~\ref{sec:practice}
adjudicates the papers the mechanical screen \emph{rejected} as well as the
ones it accepted, measures the screen's sensitivity at \auditSens, and finds
eligible papers that do report the increment across a range of operating
points --- one of them at four values of $k$ in a top-$k$ evaluation. The
earlier ``not one'' was a property of a coder that finds fewer than half of
the eligible papers, not a property of the literature. The corrected estimate
is \auditRangeThresholdPct\ with an interval that includes it.

\paragraph{Class F --- internal inconsistency}
Four corrections (C17, C18, C19, C22) are places where an earlier manuscript
disagreed with itself: an instrument count that treated two thresholds of one
instrument as two instruments; a reference reduction quoted as
$43.7\%$ in the text and $35.0\%$ in a table without the change of
specification being made clear; and a section that said a withdrawn
instrument was retained in a table that did not contain it. The generated-macro
discipline of \ref{app:verifier} makes this class structurally
impossible rather than merely unlikely --- and does not make it impossible
when the same wrong quantity is computed twice from the same wrong expression,
which is C22. Correcting the pilot's frame size in the facts and not in the
PRISMA flow left Section~\ref{sec:practice} reading
\nAuditFrame\ and Table~\ref{tab:prisma} reading two orders of magnitude more, four pages apart,
\emph{after} the error had been found and written up. The frame is computed
once now, before both consumers, and the verifier checks every row of the flow
that the prose also quotes.

\paragraph{Class G --- an estimand that was not the one being estimated}
One correction (C20), found in round nineteen and the reason
\ref{app:noisy} exists. The simulation of Section~\ref{sec:sim}
computed both estimands by enumerating the data generator's cells, which in
the noisy world carry the register's \emph{true} value while the estimator
only ever observes a value replaced by a uniform draw with probability
\simNoiseRate. Coverage was therefore measured against a population the
estimator does not sample from, and came out at \coverageNoisyBefore\ with an
apparent bias of \biasNoisyBefore. Both were artefacts. The correction is
exact marginalisation of the noise; the estimator's mean was
\meanNoisyEstimate\ before and after.

This class is worth separating from Class F because nothing in the
generated-macro discipline could have caught it: every number was internally
consistent, correctly computed, and correctly printed. What was wrong was the
\emph{definition} of the quantity two of them referred to. It was found by
constructing two experiments that could each have confirmed the comfortable
reading --- that the method simply fails in that world --- and having both
refuse to. That is the only mechanism in this project that has ever caught an
error of this class, and it is not automatable.

\paragraph{Class H --- a statistic aggregated over the wrong index set}
One correction (C21). The practice pilot's frame size summed a stratum-size
column over the \emph{rows} of the sample rather than over its distinct
strata, and a stratum's size is carried on every row of that stratum. The
quantity computed was $\sum_h N_h n_h$ and the quantity intended was
$\sum_h N_h$, so a frame of \nAuditFrame\ records was reported as one two
orders of magnitude larger. Every input was correct; the arithmetic was
correct; the group-by was wrong.

The correction changed what the pilot \emph{is} and not only how large it
sounds. With the true frame in view it is immediate that the design weight is
\auditMaxWeight\ in every stratum, that the enumeration stopped short of its
declared budget of \nAuditTarget, and that the pilot is therefore a census of
what one index returned rather than a probability sample of a literature ---
which Section~\ref{sec:practice} now says, with the additional limit that the
frame's coverage of the field is unmeasured.

The guard is a re-derivation rather than a recount: the verifier computes the
frame from the sample by grouping, and refuses a manuscript in which a frame
is smaller than the sample it contains, or differs from it while every design
weight is one. A checker that recomputed the same aggregation would have
agreed with the error.

\paragraph{Class I --- an inferential object whose denominator was not the
one the claim quantified over}
Two corrections (C26, C28), and they are the same mistake in two places. A
region label such as \emph{uniformly beneficial} ranges over every admissible
cell of a surface, and an earlier version supported it with max-$t$ bands
whose family was the five scalar instruments \emph{at one cell}; the family
was two orders of magnitude smaller than the claim. Widening it to the whole
surface raises the critical value from \maxTWithinMedian\ to
\maxTSurfaceMedian\ at the median pair and changes the region label on
\nRegionChangedByFamily\ of \nPairs\ pairs --- and more at the cell level,
where a region label's coarseness no longer hides it and the wider family
unresolves cells throughout. Section~\ref{sec:simbands} now
declares four interval objects by name so a table cannot print one and be read
as the other. The same class covers the
pilot's Wilson intervals (C28): a binomial interval was put around a ratio of
two weighted totals through a fractional ``effective sample size'', and none
of the five resulting intervals contains its design-based counterpart.

The general form is worth stating because it is not a coding error and no
checker in this repository would have caught it: \textbf{an interval is only
as good as the correspondence between the family it controls and the sentence
it is used to license}, and that correspondence is a judgement about English,
not a property of the arithmetic.

\paragraph{Class J --- a quantity assembled from units that do not add}
Two corrections (C25, C27). The first: an earlier version reported the largest
per-axis difference $\max_i(S_{T_i} - S_i)$ and described it as the share of
variance carried by interactions. Those differences overlap across axes ---
a two-way component appears in the total effect of both of its axes --- so
their maximum is neither their sum nor their union, and the total higher-order
share is $1 - \sum_i S_i$, which on this corpus is \interactionTotalPct\
against the \largestInvolvementPct\ that was printed. The second: expected
regret was averaged over an admissible set spanning five scalar instruments
and reported as a single percentage ``in the metric's own units''. There is no
such unit once AUC and Brier skill are inside one average.
Section~\ref{sec:rules} now reports regret inside one instrument at a time,
and Section~\ref{sec:nbregret} reports a common-utility version on calibrated
net benefit, which is a scale on which averaging is defined.

\paragraph{Class K --- a theorem stated wider than its proof, and wider than
the truth}
Two corrections (C23, C24). The necessity direction of the rank-invariance
proposition quantified over every non-rank-based instrument and was supported
by one stored configuration for five named instruments under one map. Repairing
the proof also refuted the claim: what invariance of the \emph{reduction}
requires is that the recalibration act \emph{affinely} on the instrument, not
that it leave it fixed, and Section~\ref{sec:props} exhibits a
non-rank-based instrument whose reduction is invariant. The word
\emph{exactly} was wrong, not merely unproved. C24 withdraws the axis
non-redundancy search entirely: its answer moved with an arbitrary tolerance,
no claim depended on it, and removing it was cheaper than repairing it.

\paragraph{Class L --- a number the design cannot produce, and a word the
design does not license}
Two corrections (C29, C30). An earlier version printed raw and Holm-adjusted
bootstrap $p$-values of $0.000$ from 200 draws; no finite bootstrap can
produce zero,
and Section~\ref{sec:planned} now uses the plus-one estimator with the
smallest attainable value stated beside it. And the five contrasts were called
\emph{confirmatory} when they had been fixed after eighteen rounds of analysis
on the same data; they are \emph{final-round planned contrasts}, which is
what the design supports.

\paragraph{Class M --- an interpretation attached to a scale that does not
carry it}
One correction (C31). Decision-curve thresholds were read as probability
thresholds,
and therefore as cost ratios, on scores whose calibration slope was reported
in the same manuscript as far from one. Section~\ref{sec:calibfirst} refits
every model contributing to a decision-curve conclusion under a training-only
calibration and labels raw-score curves as score-threshold sensitivity
analyses, which is what they are.

\paragraph{What found these, and what to take from it}
Of the corrections above, those found in the project's own work divide by
\emph{method} rather than by class. The estimand error (C20) was found by
constructing two experiments that could each have confirmed the comfortable
reading and having both refuse. The aggregation error (C21) and the
inconsistency it left behind (C22) were found by reading the compiled pages.
So were five defects that produced no wrong number at all but printed badly:
the word \emph{Appendix} doubled on every appendix reference, a macro eating
the space before an em dash, a macro carrying words set as italic variables, a
sentence's tail printed twice, and a minus sign set as a hyphen in the
sentence about signs.

Every checker in this repository ran clean on the build that carried all of
them. That is the finding worth generalising: an apparatus that guarantees
every number is derived from a result file guarantees nothing about whether
the sentence around it says what the author meant, and the only instrument
that has ever caught that is somebody reading the output.
